\scanpage{73}
% archive foliation: 72

\deleted{pour une section}
considérons (pour $S,X,t$ fixés) la catégorie des $(X',t',\underline G)$, \deleted{où $X'$ est un $X$-schéma} \inserted{tels que dessus}, comme morphismes un
\[
(X',t',G)\mathbin{-}(X'_1,t'_1,G_1)
\]
\deleted{\unsure{étant}} couple \inserted{de} $X'\mathbin{-}X'_1$ ($X$-morphisme et de $G\mathbin{-}G_1$, compatibles dans un sens évident avec les sections marquées et les opérations.
\margin{\uline{NB} On n'exige pas que $n=n_1$. Pour $X',t',G$ fixés, $(X'_1,t'_1,G_1)$ est déterminé à isom can. près, quand on se donne le quotient $G\longrightarrow G_1$.}
Considérons d'autre part les \deleted{triples} systèmes
\[
(\underline G,\unsure{\deleted{\xi}}n,\deleted{\unsure{\xi}}i,\deleted{\unsure{G}}\unsure{\xi|\zeta},T,\kappa).
\]
où \inserted{$\underline G$ un schéma en groupes fini étale sur $S$,} $n\in\mathbb N^*$, $\unsure{\xi|\zeta}\in H^0(S,R^1f^o_*f^{o*}(\underline G))$ tels que les rev. locaux (ét. loc. sur $S$)) de $X$ qu'elle décrit soient modérément ramifiés d'ordre $n$, le long de $t$, $i:\mu_{nS}\hookrightarrow\underline G$ est un mono. de groupes, $T$ un Module inv. sur $S$, et $\kappa$ un isom
\[
T^{\otimes n}\xrightarrow{\sim}T_{X/S,t}
\]
avec la définition ad hoc des morphismes (\uline{NB} \unsure{tout hom.|tout un} $\underline G\mathbin{-}\underline G_1$ doit induire un épimorphisme $\mu_{nS}\mathbin{-}\mu_{n_1S}$, qui soit l'hom \deleted{\unsure{que}} $\lambda\longmapsto\lambda^{n/n_1}$ \ldots). On a un foncteur canonique de la première catégorie dans la
